\author{Geoffrey Khan and Paul Noorlander and Katherine Hodgson and Masoud Mohammadirad and Shuan Osman Karim and Ergin \"Opengin and Eleanor Coghill and Mahîr Dogan and Murad Suleymanov and Dorota Molin and Geoffrey Haig and Parvin Mahmoudveysi} %use this field for editors as well
\title{North-Eastern Neo-Aramaic in the Context of the Languages of Western Asia}
% \title{North-Eastern Neo-Aramaic in the Linguistic Ecology of Western Asia}
% \subtitle{}
% \dedication{With contributions by Eleanor Coghill, Mahîr Dogan, Murad Suleymanov, and Dorota Molin}
% \renewcommand{\lsImpressumCitationText}{Geoffrey Khan, Paul Noorlander, Katherine Hodgson, Masoud Mohammadirad, Shuan Osman Karim and Ergin \"Opengin with contributions by Eleanor Coghill, Mahîr Dogan, Murad Suleymanov, Dorota Molin, Geoffrey Haig, and Parvin Mahmoudveysi. 2026. \textit{North-Eastern Neo-Aramaic in the Context of the Languages of Western Asia}. Berlin: Language Science Press.}
% \ISBNdigital{}
% \ISBNhardcover{}
% \BookDOI{}
% \typesetter{}
% \proofreader{}
 \lsCoverTitleSizes{48pt}{16pt}
\renewcommand{\lsSeries}{cam}
% \newcommand{\lsSeriesTitle}{}
% \newcommand{\lsSeriesColor}{}
\renewcommand{\lsSeriesNumber}{} 
\BackBody{\textit{North-Eastern Neo-Aramaic in the Context of the Languages of Western Asia} offers a comprehensive investigation of the development of North-Eastern Neo-Aramaic (NENA) through sustained contact with the languages of one of the world’s most linguistically diverse regions. Bringing together evidence from NENA and its principal contact languages—including Kurdish, Gorani, Zazaki, Armenian, Arabic and Turkic—the volume explores how centuries of multilingual interaction have shaped the phonology, morphology and syntax of these endangered Aramaic varieties. The book brings together leading experts in the respective languages and offers a detailed description of most of the internal diversity documented to date.

Drawing on extensive documentation of NENA dialects and detailed comparative analysis, the book examines a wide range of linguistic phenomena, from sound systems, nominal modification and verbal morphology to alignment, word order and argument marking. Rather than treating shared structures simply as products of contact, it investigates the historical, geographical and social conditions under which linguistic convergence has taken place. Particular attention is paid to the effects of migration, patterns of settlement and the boundaries of social and communal networks, which can produce striking differences even between communities occupying the same geographical area.

The volume also makes a broader contribution to the study of language contact. Its concluding synthesis develops a fine-grained account of contact-induced change, distinguishing processes such as matter borrowing, pattern replication, simplification, enrichment and grammaticalisation, while showing how replication may be partial, functionally transformed or shaped by multiple causes. It demonstrates that linguistic change cannot be understood through geographical proximity alone: social networks, historical layering, internal developments and speakers’ interpretations of structures in other languages all play crucial roles.

Combining rich empirical documentation with comparative and theoretical perspectives, this book will be an essential resource for scholars of Aramaic and the languages of Western Asia, as well as for researchers interested in historical linguistics, linguistic typology, multilingualism and the dynamics of language contact.}
